\subsection{Public-key cryptography from LWE}\label{lwe-crypto}

In Section \ref{lwe-small} we saw how even a small case of this problem
(\(q = 11\), \(n = 4\)) can be annoyingly tricky. In the real world, you
should imagine that \(n\) and \(q\) are much bigger -- maybe \(n\) is in
the range \(100 \leq n \leq 1000\), and \(q\) could be anywhere from
\(n^{2}\) to \(2^{\sqrt{n}}\), say.

As an example of how LWE can be used, let us see how to turn LWE into a
public-key cryptosystem. We will use the same numbers from the ``blue
set'' in Section \ref{lwe-small}. In fact, that ``blue set'' will be exactly the
public key.

\begin{center}
\begin{tabular}{|c|}
\hline
Public Key \\
\hline
\hline
(1, 0, 1, 7) : 2 \\
\hline
(5, 8, 4, 10) : 2 \\
\hline
(7, 7, 8, 5) : 3 \\
\hline
(5, 1, 10, 6) : 10 \\
\hline
(8, 0, 2, 4) : 9 \\
\hline
(9, 3, 0, 6) : 9 \\
\hline
(0, 6, 1, 6) : 9 \\
\hline
(0, 4, 9, 7) : 5 \\
\hline
(10, 7, 4, 10) : 10 \\
\hline
(5, 5, 10, 6) : 9 \\
\hline
(10, 7, 3, 1) : 9 \\
\hline
(0, 2, 5, 5) : 6 \\
\hline
(9, 10, 2, 1) : 3 \\
\hline
(3, 7, 2, 1) : 6 \\
\hline
(2, 3, 4, 5) : 3 \\
\hline
(2, 1, 6, 9) : 3 \\
\hline
\end{tabular}
\end{center}

The private key is simply the vector \(a\).

\begin{center}
\begin{tabular}{|c|}
\hline
Private Key \\
\hline
\hline
\(\mathbf{a}\) = (10, 8, 10, 10) \\
\hline
\end{tabular}
\end{center}

Since the LWE problem is hard, we can release the public key to
everybody, and they will not be able to determine the private key.

\subsubsection{Encryption}\label{how-to-encrypt-mu}

Suppose you have a message \(m \in \left\{ 0,5 \right\}\). (You will see
in a moment why we insist that \(m\) is one of these two values.) The
ciphertext to encrypt \(m\) will be a pair
\(\left( \mathbf{x}:y \right)\), where \(x\) is a vector, \(y\) is a
scalar, and \(\mathbf{x} \cdot \mathbf{a} + \epsilon = y + m\), where
\(\epsilon\) is ``small''.

How to do the encryption? If you are trying to encrypt, you only have
access to the public key -- that list of pairs
\(\left( \mathbf{x}:y \right)\) above. You want to make up your own
\(\mathbf{x}\), for which you know approximately the value
\(\mathbf{x} \cdot \mathbf{a}\). You could just take one of the vectors
\(\mathbf{x}\) from the table, but that would not be very secure: If I
see your ciphertext, I can find that \(\mathbf{x}\) in the table and use
it to decrypt \(m\).

Instead, you are going to combine several rows of the table to get your
vector \(\mathbf{x}\). Now you have to be careful: When you combine rows
of the table, the errors will add up. We are guaranteed that each row of
the table has \(\epsilon\) either \(0\) or \(1\). So if you add at most
four rows, then the total \(\epsilon\) will be at most \(4\). Since
\(m\) is either \(0\) or \(5\) (and we are working modulo \(q = 11\)),
that is just enough to determine \(m\) uniquely.

So, here is the method. You choose at random four (or fewer) rows of the
table, and add them up to get a pair \(\left( \mathbf{x}:y_{0} \right)\)
with \(\mathbf{x} \cdot \mathbf{a} \approx y_{0}\). Then you take
\(y = y_{0} - m\) (mod \(q = 11\) of course), and send the message
\(\left( \mathbf{x}:y \right)\).

\subsubsection{An example}\label{an-example}

Let us say you randomly choose the four rows:

\begin{center}
\begin{tabular}{|c|}
\hline
Some rows of public key \\
\hline
\hline
(1, 0, 1, 7) : 2 \\
\hline
(5, 8, 4, 10) : 2 \\
\hline
(7, 7, 8, 5) : 3 \\
\hline
(5, 1, 10, 6) : 10 \\
\hline
\end{tabular}
\end{center}

Now you add them up to get the following.

\begin{center}
\begin{tabular}{|c|}
\hline
\(\mathbf{x}:y_{0}\) \\
\hline \hline
(7, 5, 1, 6) : 6 \\
\hline
\end{tabular}
\end{center}

(For reference, the actual value is \(4\), so our accumulated error is
\(2\).)

Finally, let us say your message is \(m = 5\). So you set
\(y = y_{0} - m = 6 - 5 = 1\), and send the ciphertext:

\begin{center}
\begin{tabular}{|c|}
\hline
\(\mathbf{x}:y\) \\
\hline \hline
(7, 5, 1, 6) : 1 \\
\hline
\end{tabular}
\end{center}

\subsubsection{Decryption}\label{decryption}

Decryption is easy! The decryptor knows
\[\mathbf{x} \cdot \mathbf{a} + \epsilon = y + m\] where
\(0 \leq \epsilon \leq 4\).

Plugging in \(\mathbf{x}\) and \(\mathbf{a}\), the decryptor computes
\[\mathbf{x} \cdot \mathbf{a} = 4.\] Plugging in \(y = 1\), we see that
\[4 + \epsilon = 1 + m.\]

Now it is a simple ``rounding'' problem. We know that \(\epsilon\) is
small and positive, so \(1 + m\) is either \(4\) or \ldots{} a little
more. (In fact, it is one of \(4,5,6,7,8\).) On the other hand, since
\(m\) is 0 or 5, \(1 + m\) had better be 1 or 6, so the only possibility
is that \(m = 5\) (so \(1 + m = 6\)).

\subsubsection{How does this work in
general?}\label{how-does-this-work-in-general}

In practice, \(n\) and \(q\) are often much larger. Maybe \(n\) is in
the hundreds, and \(q\) could be anywhere from ``a little bigger than
\(n\)'' to ``almost exponentially large in \(n\),'' say
\(q = 2^{\sqrt{n}}\). In fact, to do FHE, we are going to want to take
\(q\) pretty big, so you should imagine that \(q \approx 2^{\sqrt{n}}\).

For security, instead of adding \(4\) rows of the public key, we want to
add at least \(\log\left( q^{n} \right) = n\log q\) rows. To be safe,
maybe a little bigger, say \(N = 2n\log q\) (of course, for this to
work, the public key has to have at least \(N\) rows). The encryption
algorithm will be ``select some subset of the rows at random, and add
them up''.

Combining \(N\) rows will have the effect of multiplying the error by
\(N\), so if the initial \(\epsilon\) was bounded by \(1\), then the
error in the ciphertext will be at most \(N\). But remember that \(q\)
is exponentially large compared to \(N\) and \(n\) anyway, so a mere
factor of \(N\) should not scare us!

To generalize our choice of \(m\) in \(\{ 0,5\}\), we could encode a
single bit by using either \(0\) or \(\lfloor\frac{q}{2}\rfloor\) to
obtain maximum separation and thus tolerance to error. Alternatively, we
could allow the message to be any multiple of some constant \(r\), where
\(r\) is bigger than the error bound, which
allows you to encode a message space of size \(q/r\) rather than just a
single bit.

When we do FHE, we are going to apply many operations to a ciphertext,
and each is going to cause the error to grow. We are going to have to
put some effort into keeping the error under control, and the size of
\(q/r\) will determine how many operations we can do before the error
grows too big.
